\begin{table}[tp]
\centering
\caption{Exact ball-proximal complexity for \href{https://github.com/k3nfalt/longtails/blob/main/LongTails/Setups/ConvexBallProximal.lean}{ConvexBallProximal}.}\label{tab:ConvexBallProximal}
\parbox{\linewidth}{\footnotesize In a finite-dimensional real inner-product space, $f$ is convex on $C$ and its epigraph restricted to $C$ is closed. A minimizer $x_\star\in C$ exists, $x_0\in C$, $\|x_0-x_\star\|\le R$, and $f(x_0)-f(x_\star)\le\Delta$, with $R\ge0$ and $t,\Delta,\varepsilon>0$. The setup fixes $x_{k+1}\in\arg\min_{z\in C,\,\|z-x_k\|\le t}f(z)$ and last-iterate output. Each step costs one exact ball-oracle call. Target: $x_K\in C$ and $f(x_K)-f(x_\star)\le\varepsilon$ for every admissible tie breaking (hence almost surely).}\par\smallskip
\begin{threeparttable}
\footnotesize
\setlength{\tabcolsep}{4pt}\renewcommand{\arraystretch}{1.25}
\begin{tabularx}{\linewidth}{@{}>{\raggedright\arraybackslash}p{0.20\linewidth}>{\raggedright\arraybackslash}p{0.32\linewidth}>{\raggedright\arraybackslash}p{0.25\linewidth}>{\raggedright\arraybackslash}X@{}}
\toprule
\textbf{Result} & \textbf{Oracle complexity} & \textbf{Novelty} & \textbf{Limitations} \\
\midrule
\textbf{(1)} \href{https://github.com/k3nfalt/longtails/blob/main/LongTails/Theorems/BallProximalSharpProfile.lean}{Ball\allowbreak Proximal\allowbreak Sharp\allowbreak Profile} \textbf{(Ours)} $(\star\star)$ \href{https://github.com/k3nfalt/longtails/blob/main/scripts/reviews/BallProximalSharpProfile.md}{(review)} & $\mathcal{O}\!\left(1+\min\!\left\{\frac{R^2}{t^2},\frac Rt h(r)\right\}\right)$ & & \\
\midrule
\textbf{(2)} \href{https://github.com/k3nfalt/longtails/blob/main/LongTails/Theorems/BallProximalGaussian.lean}{Ball\allowbreak Proximal\allowbreak Gaussian} \textbf{(Ours)} \href{https://github.com/k3nfalt/longtails/blob/main/scripts/reviews/BallProximalGaussian.md}{(review)} & $\mathcal{O}\!\left(1+\min\!\left\{\frac{R^2}{t^2},\frac Rt \sqrt\ell\right\}\right)$ & & \\
\midrule[1.2pt]
\textbf{(3)} \citep{RichtarikLi2026BallV1}, Thm.~3.4, (42)--(43)
& $\mathcal{O}\!\left(1+\min\!\left\{\frac{R^2}{t^2},\frac Rt\sqrt\ell\right\}\right)$
& \textbf{(1)}: Same order; uniformly no larger budget than (43). At $r=1/5$, $R/t=1000$: $842$ versus $897$ calls. \textbf{(2)}: Matches known rate; Lean formalization.
& \textbf{(2)}: Gaussian relaxation loses the sharper rational guarantee in (42). \\
\midrule
\textbf{(4)} \citep{RichtarikLi2026BallV1}, Rem.~3.5
& $\mathcal{O}\!\left(1+\frac Rt\sqrt{(1-r^2)\log(1/r)}\right)$\newline Fixed $0<r<1$, $R/t\to\infty$; implicit finite-step bound.
& \textbf{(1)}: Sharp leading constant via supporting worst-case characterization. At $r=1/5$, $R/t=1000$: $842$ versus $\approx879$ calls; gain vanishes as $r\to0$.
& \textbf{(1)}: No finite-step dominance; \textbf{(4)} is exact at $K=1$. \textbf{(2)}: Weaker gap guarantee. \\
\midrule
\textbf{(5)} \citep{RichtarikLi2026BallV1}, Thm.~4.5
& $\mathcal{O}\!\left(1+\min\!\left\{\frac{R^2}{t^2},C_d\frac Rt\right\}\right)$\newline Exact minimizer.
& \textbf{(1)}, \textbf{(2)}: Lower target-accuracy cost when $1\ll\sqrt\ell\ll\min\{C_d,R/t\}$: $(R/t)\sqrt\ell$ versus $\min\{(R/t)^2,C_dR/t\}$.
& \textbf{(1)}, \textbf{(2)}: Worse when $C_d\ll\min\{\sqrt\ell,R/t\}$; \textbf{(5)} attains the exact optimum. \\
\midrule
\textbf{(6)} \citep{GruntkowskaEtAl2025BallV2}, Thm.~8.1(iii), Cor.~8.2
& $\mathcal{O}\!\left(1+\min\!\left\{\frac{R^2}{t^2},\frac{\ell}{\log(1+t/R)}\right\}\right)$\newline $R>0$.
& \textbf{(1)}, \textbf{(2)}: $\sqrt\ell$ instead of $\ell$ when $1\ll\ell\ll R/t$; total saving $\Theta(\sqrt\ell)$, already achieved by \textbf{(3)}.
& \\
\bottomrule
\end{tabularx}
\begin{tablenotes}[flushleft]\footnotesize
\item[] \textbf{Abbr:} $r=\varepsilon/\Delta$, $\ell=\log_+(\Delta/\varepsilon)$, $\log_+u=\max\{0,\log u\}$; $h(r)=(1-r)_+$ for $r\ge1/2$ and $h(r)=\sqrt{\log(1/r)}$ for $0<r<1/2$; $(u)_+=\max\{0,u\}$. $d=\dim E$; $C_d$ is the dimension-dependent self-contracted-curve length constant ($C_1=2$, $C_d=4\sqrt d(1+4\sqrt d)^d$ for $d\ge2$). Costs count exact ball-oracle calls; $\log$ is natural. Numerical comparisons retain the original budgets' constants.
\end{tablenotes}
\end{threeparttable}
\end{table}
\clearpage
